\problem{}
\medskip

\textbf{Monotone 3-SAT:} A \emph{3-SAT formula} is a Boolean formula in which each clause contains exactly three literals. A clause is called \emph{monotone} if all of its literals are positive or all of its literals are negative. For example, $(x \vee y \vee z)$ and $(\neg x \vee \neg y \vee \neg z)$ are monotone clauses, while $(x \vee y \vee \neg z)$ is not. As usual, a clause is satisfied if at least one of its literals evaluates to true.

Given a 3-SAT formula $\phi$ in which every clause of $\phi$ is monotone, decide whether $\phi$ has a satisfying assignment. Prove that \textsc{Monotone 3-SAT} is NP-complete.
\medskip


\solution{
\textbf{1. Membership in NP.}
The problem is in NP, since a truth assignment can be checked by evaluating all
clauses.

\textbf{2. Known NP-complete problem.}
We reduce from \textsc{3-SAT}, which is NP-complete.

\textbf{3. Polynomial-time reduction.}
Let $\phi$ be an instance of \textsc{3-SAT}, where every clause has exactly
three literals. We transform each clause separately. Clauses that are already
monotone are left unchanged. We use the usual convention that the three
literals are three literal occurrences; the statement does not require them to
be three distinct variables.

For a mixed clause with two positive literals and one negative literal, write it
as
\[
(x\vee y\vee \neg z).
\]
Introduce a fresh variable $u$ for this clause and replace the clause by
\[
(x\vee y\vee u)\ \wedge\ (\neg u\vee \neg z\vee \neg z).
\]
Both clauses are monotone.

For a mixed clause with one positive literal and two negative literals, write it
as
\[
(x\vee \neg y\vee \neg z).
\]
Introduce a fresh variable $u$ for this clause and replace the clause by
\[
(x\vee u\vee u)\ \wedge\ (\neg u\vee \neg y\vee \neg z).
\]
Again both clauses are monotone. The construction adds at most one fresh
variable and two clauses for each original clause, so the size only increases
linearly.

\textbf{4. Correctness of the reduction.}
Now we check that each replacement preserves satisfiability.

Consider first a clause
\[
C=(x\vee y\vee \neg z)
\]
and its replacement
\[
C'=(x\vee y\vee u)\wedge(\neg u\vee \neg z\vee \neg z).
\]
If $C$ is true and at least one of $x,y$ is true, set $u=\mathrm{false}$. If
both $x,y$ are false, then $z$ must be false, and we set $u=\mathrm{true}$.
In both cases $C'$ is true. Conversely, if $C'$ is true, then either $x$ or $y$
is true, in which case $C$ is true, or both $x,y$ are false. In the latter case
the first clause forces $u=\mathrm{true}$, and then the second clause forces
$z=\mathrm{false}$. Thus $C$ is also true.

The other mixed case is symmetric. For
\[
C=(x\vee \neg y\vee \neg z)
\]
with replacement
\[
C'=(x\vee u\vee u)\wedge(\neg u\vee \neg y\vee \neg z),
\]
if $C$ is true, set $u=\mathrm{false}$ when $x$ is true, and set
$u=\mathrm{true}$ otherwise. Conversely, if $C'$ is true and $x$ is false, then
the first clause forces $u=\mathrm{true}$, so the second clause forces at least
one of $\neg y,\neg z$ to be true. Hence $C$ is true.

So an assignment satisfying $\phi$ can be extended to the new variables, and any
satisfying assignment of the new formula restricts to a satisfying assignment of
$\phi$. Therefore
\[
\phi\text{ is satisfiable}
\quad\Longleftrightarrow\quad
\phi'\text{ is satisfiable}.
\]
Hence \textsc{Monotone 3-SAT} is NP-hard. It is also in NP, so it is
NP-complete.
}

\newpage
